\section{Engineering Requirements Specification}
\label{sec:requirements}

\subsection{Formal Requirements Classification under VDI 2206}
The task clarification and conceptual design of complex mechatronic systems must proceed from a structured, unambiguous specification of engineering requirements. In accordance with the VDI~2206 design guideline (\emph{Entwicklungsmethodik f{\"u}r mechatronische und cyber-physische Systeme})~\cite{vdi2206}, qualitative user expectations, operational mission profiles, and regulatory constraints are systematically converted into a formal engineering requirements list (\emph{Anforderungsliste}). This document serves as the inviolable baseline against which all downstream architectural, mechanical, electrical, and algorithmic decisions are evaluated and validated across the macro-cycles and micro-cycles of the V-model.

To enforce engineering rigor and avoid premature design convergence, each requirement is categorized according to two formal classification axes:
\begin{enumerate}[leftmargin=1.5em]
  \item \textbf{Requirement Severity Type:}
  \begin{itemize}[leftmargin=1.2em]
    \item \textbf{Fixed Requirement (\emph{Festforderung}, FF / Demand):} Inviolable boundary constraints that must be fulfilled unconditionally under all nominal and degraded operating conditions. An architectural concept or hardware implementation that violates a single fixed requirement is technically disqualified from consideration.
    \item \textbf{Range Requirement (\emph{Bereichsforderung}, BF / Wish):} Performance criteria that specify target bands with defined acceptable minimal baselines and optimal target goals. The engineering embodiment must satisfy the lower bound, while design optimizations and trade-off studies seek to maximize performance toward the upper target.
  \end{itemize}
  \item \textbf{Criterion Verification Characteristic:}
  \begin{itemize}[leftmargin=1.2em]
    \item \textbf{Quantitative (\emph{quan}):} Physical and operational parameters defined by numerical metrics with standard SI engineering units. These criteria are empirically verifiable via finite element analysis (FEA), multi-body dynamic simulation, coordinate measuring, calibrated physical instruments (e.g., laser trackers, dynamometers, oscilloscopes), or sensor telemetry.
    \item \textbf{Qualitative (\emph{qual}):} Functional, topological, architectural, or ergonomic properties that cannot be expressed as a continuous scalar metric. Verification is performed through formal design audits, architectural code compliance reviews, protocol packet analyzers, interface compliance testing, or safety certification checklists.
  \end{itemize}
\end{enumerate}

By anchoring these constraints to the standardized operational benchmarks of domestic and service robotics---most notably the RoboCup@Home General Purpose Service Robot (GPSR), Table Clearing, and Articulated Fixture challenges~\cite{robocupathome}---the resulting specification directly reflects the physical realities of human-centric indoor spaces.

\subsection{Comprehensive Engineering Requirements List (\emph{Anforderungsliste})}
Table~\ref{tab:anforderungsliste} details the complete engineering requirements specification governing the modular semi-humanoid robot platform across geometry, kinematics, dynamics, electrical power, embedded control, sensing, software, human-machine interface (HMI), and collaborative safety.

\begin{small}
\setlength{\tabcolsep}{4pt}
\begin{longtable}{@{}p{0.06\textwidth}p{0.19\textwidth}p{0.25\textwidth}p{0.06\textwidth}p{0.06\textwidth}p{0.25\textwidth}@{}}
\caption{Comprehensive engineering requirements list (\emph{Anforderungsliste}) for the modular semi-humanoid robot platform in accordance with VDI 2206.}
\label{tab:anforderungsliste}\\
\toprule
\textbf{ID} & \textbf{Requirement Parameter} & \textbf{Target Specification} & \textbf{Type} & \textbf{Char.} & \textbf{Verification Method} \\
\midrule
\endfirsthead
\multicolumn{6}{c}{{\bfseries Table \thetable\ continued from previous page}} \\
\toprule
\textbf{ID} & \textbf{Requirement Parameter} & \textbf{Target Specification} & \textbf{Type} & \textbf{Char.} & \textbf{Verification Method} \\
\midrule
\endhead
\midrule
\multicolumn{6}{r}{{Continued on next page\dots}} \\
\endfoot
\bottomrule
\endlastfoot

\multicolumn{6}{@{}l}{\textbf{1. Geometry and Mass (\emph{Geometrie und Masse})}} \\
\midrule
G01 & Base Footprint Envelope & Diameter / width envelope $\le 550\text{ mm}$ & BF & quan & Dimensional CAD audit \& chassis jig check \\
G02 & Doorway Passage Clearance & Passage through doorways $\ge 750\text{ mm}$ (clearance margin $\ge 100\text{ mm}$/side) & FF & quan & Physical passage trial in benchmark doorway arena \\
G03 & Overall Platform Height & Standing height $1200\text{--}1300\text{ mm}$ (nominal home configuration) & BF & quan & Vertical height measurement from floor plane \\
G04 & Total System Mass & Total wet mass $\le 55\text{ kg}$ (including battery \& all payloads) & FF & quan & Calibrated multi-point platform scale measurement \\
\midrule

\multicolumn{6}{@{}l}{\textbf{2. Kinematics and Workspace (\emph{Kinematik und Arbeitsraum})}} \\
\midrule
K01 & Total Active Degrees of Freedom & Total active $\text{DoF} = 18$ (Base: 2, Arms: $2 \times 6 = 12$, Grippers: 2, Head: 2) & FF & quan & Joint topological audit \& actuator allocation check \\
K02 & Mobile Base Kinematic Drive & Differential drive (2 active motorized wheels + 2 passive spring-loaded casters) & FF & qual & Kinematic dead-reckoning test \& zero-radius spin \\
K03 & Upper-Body Arm Topology & Dual 6-DoF articulated serial chains (anthropomorphic shoulder-elbow-wrist) & FF & qual & Forward/inverse kinematic Jacobian rank audit \\
K04 & Vertical Manipulation Reach & Continuous vertical reach envelope $500\text{--}1350\text{ mm}$ above finished floor level & FF & quan & 3D coordinate measuring of end-effector flange \\
K05 & Horizontal Bimanual Span & Maximum bimanual arm tip-to-tip span $\ge 1200\text{ mm}$ ($1200\text{--}1400\text{ mm}$) & BF & quan & Cartesian boundary mapping in simulation \& rig \\
K06 & Active Head Perception Unit & 2-DoF pan-tilt neck (Pan: $\pm 90^\circ$, Tilt: $+45^\circ / -30^\circ$) & BF & quan & Digital goniometer joint angle sweep verification \\
\midrule

\multicolumn{6}{@{}l}{\textbf{3. Dynamics and Manipulation Performance (\emph{Dynamik und Manipulation})}} \\
\midrule
D01 & Manipulator Payload Capacity & Continuous payload $\ge 1.0\text{ kg}$/arm nominal (peak/holding $\ge 1.5\text{ kg}$) & FF & quan & Static load-holding dynamometer test at reach \\
D02 & End-Effector Clamping & Dual 2-finger parallel stroke $0\text{--}85\text{ mm}$, clamping force $15\text{--}35\text{ N}$ & FF & quan & Calibrated load cell pinch test with silicone pads \\
D03 & Maximum Base Transit Speed & Regulated indoor linear velocity $v_{\text{base}} = 0.6\text{--}0.8\text{ m/s}$ on level floor & BF & quan & Laser-timed transit; stopping distance $\le 0.4\text{ m}$ \\
D04 & Cartesian TCP Repeatability & Unidirectional repeatability $\le \pm 2.0\text{ mm}$ under nominal $1.0\text{ kg}$ payload & BF & quan & Multi-cycle dial-indicator / laser tracker test \\
D05 & Settling Time \& Damping & Re-positioning settling time $t_{\text{settling}} \le 0.5\text{ s}$ (damping ratio $\zeta \ge 0.6$) & BF & quan & Flange accelerometer vibration decay recording \\
\midrule

\multicolumn{6}{@{}l}{\textbf{4. Electrical and Power Infrastructure (\emph{Elektrische Energieversorgung})}} \\
\midrule
E01 & System DC Bus Voltage & Nominal $24\text{ V DC}$ bus ($8\text{S LiFePO}_4$, operational band $20.0\text{--}28.8\text{ V}$) & FF & quan & Oscilloscope \& multi-meter bus voltage logging \\
E02 & Operational Run-Time & Run-time $\ge 3.0\text{ h}$ under continuous mixed operational duty cycle & BF & quan & Automated cyclic mobile manipulation test \\
E03 & Emergency Disconnect Latency & Hardware E-stop disconnect response time $\le 10\text{ ms}$ on actuator power rails & FF & quan & Oscilloscope capture of contactor de-energization \\
E04 & Power Rail Isolation & Galvanically isolated logic rails ($19\text{V}$ GPU, $12\text{V}$ sensors, $5\text{V}/3.3\text{V}$ logic) & FF & qual & Common-mode noise FFT analysis ($>40\text{ dB}$ reject) \\
\midrule

\multicolumn{6}{@{}l}{\textbf{5. Embedded Control and Communication (\emph{Steuerung und Kommunikation})}} \\
\midrule
C01 & Real-Time Joint Control Rate & Joint-level closed-loop frequency $\ge 500\text{ Hz}$ (target: $1000\text{ Hz}$, jitter $\le 50\ \mu\text{s}$) & FF & quan & `PREEMPT\_RT` kernel cyclic latency benchmark \\
C02 & Fieldbus Communication Bus & CAN FD / CAN 2.0B fieldbus ($\ge 1\text{ Mbps}$, bus load $\le 65\%$) & FF & qual & Bus analyzer packet inspection \& CRC error check \\
C03 & Device-Level Abstraction & Modular device abstraction for base, arm, neck, and gripper subsystems & FF & qual & Standalone sub-assembly hardware-in-the-loop test \\
\midrule

\multicolumn{6}{@{}l}{\textbf{6. Sensing and Perception (\emph{Sensorik und Umgebungserfassung})}} \\
\midrule
S01 & Base Obstacle Sensing & $360^\circ$ planar 2D LiDAR coverage + ring of ultrasonic transducers & FF & quan & Blind-spot mapping with standard reflection targets \\
S02 & 3D Visual Perception & Head-mounted stereo RGB-D sensor (depth range $0.2\text{--}3.0\text{ m}$, $\ge 720\text{p}@30\text{fps}$) & FF & qual & Point cloud calibration \& depth RMS error check \\
S03 & Contact \& Effort Feedback & Motor phase current sensing ($\pm 50\text{ mA}$ resolution) + compliant finger pads & BF & qual & Sensorless contact detection \& slip limit test \\
\midrule

\multicolumn{6}{@{}l}{\textbf{7. Software and Embodied Intelligence (\emph{Software und KI})}} \\
\midrule
A01 & Standard Middleware & ROS 2 LTS distribution (Humble Hawksbill / Jazzy Jalisco) & FF & qual & Node graph audit \& package build verification \\
A02 & Decoupled Task Port Actions & Standard Action interfaces (\texttt{/navigate}, \texttt{/pick\_object}, \texttt{/place\_object}) & FF & qual & Automated integration test flashing Task Port client \\
A03 & Spoken Command Latency & Whisper Speech-to-Text inference latency $\le 1.5\text{ s}$ for 5-second audio input & BF & quan & Onboard edge compute benchmark execution \\
A04 & VLA Action Bandwidth & Vision-Language-Action policy proposal rate $\ge 10\text{ Hz}$ with $1\text{ kHz}$ spline & BF & quan & Continuous end-to-end action streaming latency \\
\midrule

\multicolumn{6}{@{}l}{\textbf{8. Human-Machine Interface and Teleoperation (\emph{HMI und Teleoperation})}} \\
\midrule
H01 & Web Teleoperation Latency & Bidirectional web dashboard teleoperation latency $\le 100\text{ ms}$ over local Wi-Fi & BF & quan & Round-trip ping \& video pipeline glass-to-glass test \\
H02 & Telemetry \& Alarms & Real-time visual telemetry, battery SoC display, and visual/audible alarms & FF & qual & Fault injection test for temperature/voltage alarms \\
H03 & Task Programming Macro & Interactive teach-in and macro recording of waypoints \& gripper events & BF & qual & Non-expert operator usability and execution test \\
H04 & Physical Bulkhead Panel & Maintenance panel with hardware E-stop, charging port, GbE, USB 3.0 & FF & qual & Physical layout inspection \& port accessibility \\
\midrule

\multicolumn{6}{@{}l}{\textbf{9. Collaborative Safety and Compliance (\emph{Sicherheit und Normen})}} \\
\midrule
SF01 & Collaborative Robot Safety & ISO 13482 compliance: speed/separation monitoring, contact force $\le 150\text{ N}$ & FF & qual & Safety zone deceleration \& bio-fidelic impact test \\
\bottomrule
\end{longtable}
\end{small}

\subsection{Engineering Rationale for the Core Architectural Decisions}
\label{subsec:engineering_rationale}
The synthesis of the engineering requirements list in Table~\ref{tab:anforderungsliste} reflects deliberate trade-offs established to maximize physical robustness, control determinism, operational reliability, and ease of maintenance within domestic and light-industrial environments. Four foundational architectural decisions anchor the entire mechatronic design:

\subsubsection{1. Differential Drive Base vs. Omnidirectional Mecanum Drive}
While omnidirectional Mecanum wheel arrangements provide three unconstrained planar degrees of freedom ($\dot{x}, \dot{y}, \dot{\theta}$), an in-depth tribological and dynamic analysis reveals severe disadvantages that violate core operational requirements in unstructured human environments:
\begin{itemize}[leftmargin=1.5em]
  \item \textbf{Traction Loss and Dead-Reckoning Drift:} Mecanum kinematics rely on passive peripheral rollers angled at $45^\circ$. When traversing mixed indoor floorings---such as low-pile carpeting, polished tile, linoleum, or thresholds with micro-contaminants---the passive rollers experience unpredictable intermittent slip. Empirical testing of service robots reveals that Mecanum bases exhibit odometric dead-reckoning errors exceeding $15\text{--}25\%$ over modest travel distances, severely degrading Kalman-filtered wheel odometry and forcing heavy reliance on scan matching. In contrast, a differential drive base employing two high-traction polyurethane drive wheels paired with two passive, spring-damped swivel casters provides non-slip rolling contact. This ensures highly deterministic wheel odometry with dead-reckoning drift well below $2\%$.
  \item \textbf{Vibrational Excitation of Structural Linkages:} The discrete geometric transitions between adjacent rollers on a Mecanum wheel introduce high-frequency periodic impact shocks during rolling. These impacts excite structural natural frequencies in the robot's upper body and neck linkages, generating image blur in the head-mounted RGB-D camera and destabilizing sensitive optical perception pipelines. The continuous rolling profile of solid differential drive wheels eliminates this vibrational excitation mode entirely.
  \item \textbf{Energetic Efficiency and Mechanical Simplicity:} Mecanum drives require four independent high-power geared motor units, whereas differential drive requires only two motorized hubs. The parasitic scrubbing friction inherent to Mecanum roller vectors consumes $30\text{--}40\%$ more electrical power for equivalent forward thrust. By selecting differential drive, the platform saves approximately $6\text{--}8\text{ kg}$ of actuator mass and significantly extends operational battery runtime beyond $3.0\text{ h}$ (E02). Zero-radius in-place rotation ($\omega_z$) provides ample maneuverability through standard $750\text{ mm}$ doorways (G02).
\end{itemize}

\subsubsection{2. Dual 6-DoF Manipulators vs. 7-DoF Redundant Kinematic Chains}
Anthropomorphic mobile manipulation literature frequently advocates for 7-DoF arms due to the kinematic redundancy of the elbow swivel angle. However, an analysis of the specific manipulation tasks required for table clearing and pick-and-place reveals that 6-DoF serial manipulators offer decisive engineering advantages:
\begin{itemize}[leftmargin=1.5em]
  \item \textbf{Sufficiency for Task-Space Reachability:} The target operational tasks require reaching positions and orientations within the special Euclidean group $SE(3) = \mathbb{R}^3 \times SO(3)$. A 6-DoF manipulator possessing an anthropomorphic shoulder-elbow-spherical wrist topology provides the exact mathematical minimum number of joints required to position and orient the end-effector arbitrarily within its reachable volume. Because the mobile base itself provides planar degrees of freedom ($x, y, \theta$), the combined mobile manipulation system already possesses $2 + 6 = 8$ degrees of freedom per arm, providing external kinematic redundancy without inflating joint complexity.
  \item \textbf{Mass, Cantilever Inertia, and Dynamic Bandwidth:} Incorporating a 7th active joint adds approximately $1.5\text{--}2.5\text{ kg}$ of concentrated mass into the upper limb. In accordance with the fundamental relation governing natural frequency:
  \begin{equation}
    \omega_n = \sqrt{\frac{K}{M}}
  \end{equation}
  adding distal joint mass dramatically increases the cantilever moment arm exerted on the shoulder base:
  \begin{equation}
    M_b = m_{\text{arm}} g L_{\text{CoM}} + m_{\text{payload}} g L_{\text{reach}}
  \end{equation}
  This lowers the structural natural frequency, which in turn inflates the mechanical settling time:
  \begin{equation}
    t_{\text{settling}} \approx \frac{4}{\zeta \omega_n}
  \end{equation}
  By eliminating the 7th joint, link mass is reduced, shifting structural resonance frequencies above $15\text{ Hz}$ and allowing settling times $t_{\text{settling}} \le 0.5\text{ s}$ ($\zeta \ge 0.6$) under high-speed point-to-point trajectory profiles (D05).
  \item \textbf{Computational Determinism and Fieldbus Bandwidth:} Resolving the kinematic inversion of 7-DoF redundant arms requires numerical pseudo-inverse damping ($\mathbf{J}^\dagger$) or secondary null-space projection tasks. These iterative solvers introduce algorithmic non-determinism, computational latency jitter, and unpredictable joint reconfigurations that complicate real-time closed-loop control. With dual 6-DoF arms, the total joint count across both limbs is fixed at 12 DoF, halving inverse kinematic computation overhead and fitting within the deterministic packet constraints of a 1~Mbps CAN fieldbus (C02).
\end{itemize}

\subsubsection{3. Fixed Ergonomic Torso vs. Actively Elevating Prismatic Lift Column}
Many research mobile manipulators incorporate a telescoping vertical column to extend vertical reach from floor level to upper cabinetry. While versatile, prismatic columns introduce structural liabilities that compromise reliability in practical deployment:
\begin{itemize}[leftmargin=1.5em]
  \item \textbf{Mechanical Binding under Cantilever Bending Moments:} A vertical prismatic lift column requires precision linear guide rails, recirculating ball bearings, and a ball-screw or timing-belt transmission. When both 6-DoF arms are extended horizontally with payloads, the resulting forward cantilever moment exerts extreme normal reaction forces across the linear bearing carriages:
  \begin{equation}
    F_{\text{reaction}} = \frac{M_{\text{cantilever}}}{d_{\text{bearing}}} = \frac{\sum m_i g L_i}{d_{\text{bearing}}}
  \end{equation}
  where $d_{\text{bearing}}$ is the vertical spacing between bearing blocks. In practice, this high moment causes mechanical binding, stick-slip vibration, accelerated wear, and potential catastrophic jamming of the lift mechanism.
  \item \textbf{Parasitic Dead Weight and Center-of-Mass Elevation:} A motorized prismatic stage, including guide rails, structural telescoping sleeves, ball screw, and brake motor, adds $10\text{--}15\text{ kg}$ of parasitic dead weight. Elevating this mass raises the system center of mass ($z_{\text{CoM}}$), drastically increasing the risk of dynamic tip-over during abrupt emergency braking maneuvers.
  \item \textbf{Targeted Functional Workspace Coverage:} An analysis of domestic and service tasks shows that the overwhelming majority of high-value manipulation targets reside within an ergonomic height band: dining tables ($720\text{--}760\text{ mm}$), kitchen countertops ($850\text{--}920\text{ mm}$), desks ($750\text{ mm}$), and eye-level storage shelves ($1100\text{--}1350\text{ mm}$). By fixing the torso height at an ergonomic shoulder center of $\approx 950\text{ mm}$ (overall robot height $1200\text{--}1300\text{ mm}$), the dual 6-DoF arms cover a continuous vertical manipulation band from $500\text{ mm}$ to $1350\text{ mm}$ (K04) purely through arm articulation. This eliminates moving torso pinch hazards, maximizes torsional and bending rigidity ($K$), lowers the center of mass into the base chassis ($z_{\text{CoM}} \le 350\text{ mm}$), and saves substantial fabrication cost.
\end{itemize}

\subsubsection{4. Dual 2-Finger Parallel Grippers vs. Multi-Fingered Anthropomorphic Hands}
Anthropomorphic five-finger hands provide high aesthetic biomimicry but suffer from severe physical and algorithmic limitations in real-world manipulation:
\begin{itemize}[leftmargin=1.5em]
  \item \textbf{Mechanical Fragility and High Maintenance:} Multi-fingered hands require between 16 and 24 active or underactuated micro-joints driven by miniature brushed DC motors, ultra-thin tendon cables, and micro-pulleys. Under continuous pick-and-place cycles, tendon cables experience severe fatigue stretching, unseating, and rupture, resulting in prohibitive maintenance downtime and replacement costs.
  \item \textbf{Grasp Planning Dimensionality and Execution Failure:} Generating stable force-closure grasps for an articulated hand requires solving high-dimensional non-linear contact optimization problems over hundreds of candidate grasp topologies. In cluttered tabletop environments, slight object pose estimation errors often lead to fingertip collisions, dropped objects, and grasp failure.
  \item \textbf{Deterministic Performance with Compliant Silicone Pads:} In contrast, a 2-finger parallel gripper driven by a single high-efficiency lead screw or rack-and-pinion transmission delivers deterministic, backlash-free bidirectional travel ($0\text{--}85\text{ mm}$ stroke) and robust clamping forces of $15\text{--}35\text{ N}$ (D02). By outfitting the rigid aluminum finger tips with molded compliant silicone pads ($\mu_s \ge 0.8$), the contact surface passively deforms around diverse geometric profiles (cylindrical bottles, rectangular boxes, spherical fruit, delicate tableware), converting complex geometric grasping into a reliable, single-parameter linear closure. Motor phase current monitoring provides direct, sensorless grasping force regulation ($\pm 50\text{ mA}$) to securely grip delicate objects without damage.
\end{itemize}

\subsection{Decoupling Capability from Intent: The Architectural Foundation}
\label{subsec:decoupling_intent}
A foundational engineering axiom of the platform is the rigorous, architectural decoupling of \emph{physical capability} from \emph{operational intent}:
\begin{quote}
\emph{The core mechatronic platform must encapsulate all deterministic physical capabilities---motion control, trajectory generation, kinematic safety limits, and low-level sensor acquisition---behind a hardened abstraction boundary, completely independent from high-level cognitive policies, speech parsers, or probabilistic AI models.}
\end{quote}
This decoupling is implemented through the \textbf{Task Port} abstraction layer. Downstream application software, external mission orchestrators, or modern Vision-Language-Action (VLA) foundation models interact with the robot solely by submitting high-level goal actions (e.g., \texttt{/navigate}, \texttt{/pick\_object}, \texttt{/place\_object}). The platform validates these goal primitives against internal kinematic boundaries, collision models, and ISO~13482 safety limits before executing them deterministically across the 1~kHz real-time fieldbus. 

This strict separation ensures that high-level software bugs, out-of-distribution neural network hallucinations, or network disconnects cannot cause physical link collisions, actuator over-torquing, or human injury, providing an extensible and safe foundation for physical AI research.
